\section{Discussion and Limitations}
The strongest evidence concerns continued semantic feedback: 33/39 eligible completions versus 6/39 for one initial plan with the same motor controllers. This effect belongs to the full protocol, which supplies later observations and more inference compute. It does not isolate fresh images at equal compute.

Learned slowing with joint correction has the highest observed waiting-study completion, but the three registered adjacent contrasts remain unresolved. It also has more later-pending drops than learned slowing alone. Neither teacher-action fit nor aggregate drop reduction certifies uniformly better contact stability. Explicit controller-goal conditioning provides no demonstrated benefit, and the execution gate that fails development admission is not deployed. The Astra--fine-tuned difference remains unresolved on both objects.

The completed visual intervention covers one cube and one reconstructed permutation. Cube perception renders a Gaussian object with hand occlusion; dice instead uses noisy simulator pose and separate motor weights. Neither establishes real-camera robustness, zero-shot object transfer or a universal controller. Complete-system baselines differ in controller access and inspection schedules. Hosted model revisions limit exact replication. Further objects, independent predictor measurements and hardware trials remain open.

\section{Conclusion}
We presented semantic in-hand reorientation with continuous local control around asynchronous visual decisions. The experiments support continued feedback and expose visual-binding and contact-control limitations. Waiting-control gains remain uncertain, and negative goal and execution ablations are retained. Object-specific dice results extend coverage while identifying startup drops as a remaining weakness.
